% Options for packages loaded elsewhere
\PassOptionsToPackage{unicode}{hyperref}
\PassOptionsToPackage{hyphens}{url}
\PassOptionsToPackage{dvipsnames,svgnames,x11names}{xcolor}
\documentclass[
  10pt,
  fontset=fandol]{ctexart}
\usepackage{xcolor}
\usepackage[margin=16mm]{geometry}
\usepackage{amsmath,amssymb}
\setcounter{secnumdepth}{-\maxdimen} % remove section numbering
\usepackage{iftex}
\ifPDFTeX
  \usepackage[T1]{fontenc}
  \usepackage[utf8]{inputenc}
  \usepackage{textcomp} % provide euro and other symbols
\else % if luatex or xetex
  \usepackage{unicode-math} % this also loads fontspec
  \defaultfontfeatures{Scale=MatchLowercase}
  \defaultfontfeatures[\rmfamily]{Ligatures=TeX,Scale=1}
\fi
\usepackage{lmodern}
\ifPDFTeX\else
  % xetex/luatex font selection
\fi
% Use upquote if available, for straight quotes in verbatim environments
\IfFileExists{upquote.sty}{\usepackage{upquote}}{}
\IfFileExists{microtype.sty}{% use microtype if available
  \usepackage[]{microtype}
  \UseMicrotypeSet[protrusion]{basicmath} % disable protrusion for tt fonts
}{}
\makeatletter
\@ifundefined{KOMAClassName}{% if non-KOMA class
  \IfFileExists{parskip.sty}{%
    \usepackage{parskip}
  }{% else
    \setlength{\parindent}{0pt}
    \setlength{\parskip}{6pt plus 2pt minus 1pt}}
}{% if KOMA class
  \KOMAoptions{parskip=half}}
\makeatother
\setlength{\emergencystretch}{3em} % prevent overfull lines
\providecommand{\tightlist}{%
  \setlength{\itemsep}{0pt}\setlength{\parskip}{0pt}}
\usepackage{amsmath,amssymb,mathtools,needspace}
\xeCJKDeclareCharClass{CJK}{"2032,"0400->"04FF,"2160->"216B,"2460->"2473,"25A0,"25C7,"25CB,"25CF}
\IfFontExistsTF{Arial Unicode MS}{\xeCJKsetup{AutoFallBack=true}\setCJKfallbackfamilyfont{\CJKrmdefault}{Arial Unicode MS}}{}
\setcounter{secnumdepth}{0}
\setlength{\emergencystretch}{2em}
\allowdisplaybreaks
\usepackage{bookmark}
\IfFileExists{xurl.sty}{\usepackage{xurl}}{} % add URL line breaks if available
\urlstyle{same}
\hypersetup{
  pdftitle={缩减技术的设计机理},
  colorlinks=true,
  linkcolor={Maroon},
  filecolor={Maroon},
  citecolor={Blue},
  urlcolor={Blue},
  pdfcreator={LaTeX via pandoc}}

\title{缩减技术的设计机理}
\author{}
\date{}

\begin{document}
\maketitle

\subsubsection{13.2.2 缩减技术的设计机理}\label{ux7f29ux51cfux6280ux672fux7684ux8bbeux8ba1ux673aux7406}

许多数值计算问题可以引进某个实数------所谓问题的\textbf{规模}来刻画其``大小''，而问题的解则是其规模为足够小，譬如规模为 1 或 0 的退化情形。求解这类问题，一种行之有效的办法是通过某种简单的运算手续逐步缩减问题的规模，直到加工得出所求的解。算法设计的这种技术称作\textbf{规模缩减技术}，简称\textbf{缩减技术}。

\textbf{缩减技术适用的一类问题是，求解这类问题的困难在于它的规模（适当定义）比较大。针对这类问题运用缩减技术，就是设法逐步缩减计算问题的规模，直到规模变得足够小时直接生成或方便地求出问题的解。}

缩减技术的设计机理可用\textbf{``大事化小，小事化了''}这句俗话来概括。

\textbf{所谓``大事化小''意即逐步压缩问题的规模。}在运用缩减技术时，``大事''是如何``化小''的呢？这个处理过程具有如下两项基本特征。

（1）\textbf{结构递归}：``大事化小''是逐步完成的，其每一步将所考察的计算模型加工成同样类型的计算模型，因而这类算法具有明晰的递归结构。

（2）\textbf{规模递减}：每一步加工前后的计算模型虽然从属于同一类型，但其规模已被压缩了。压缩系数愈小，算法的效率愈高。

再考察``小事化了''的处理过程。\textbf{所谓``小事化了''，是指当问题的规模变得足够小时即可直接或方便地得出问题的解。}

\href{../../source-images/pdf-310.jpeg}{原页图像}

``小事''是如何``化了''的呢？

对于某些计算模型，如前面讨论过的数列求和问题，它们的规模为正整数，而其解则是规模为 0 或 1 的退化情形。这时只要设法使规模逐次减 1，加工若干步后即可直接得出所求的解。这里``小事化了''是直截了当的。

这样设计出的一类算法统称\textbf{直接法}。前述数列求和的累加算法以及下面将要讲到的多项式求值的秦九韶算法都是直接法。

\end{document}
